% Part: counterfactuals
% Chapter: introduction
% Section: material-conditional

\documentclass[../../../include/open-logic-section]{subfiles}

\begin{document}

\olfileid{cnt}{int}{mat}

\olsection{实质条件句}

英语中最简单的条件句形式是“如果……那么……”这样的语句，其中“……”本身也是语句，例如“如果管家是凶手，那么园丁就是无辜的”。在入门逻辑课程中，我们学会用 $\lif$ 联结词来符号化条件句：将“……”所指示的部分（例如，用公式 $!A$ 和 $!B$）符号化，整个条件句就被符号化为 $!A \lif !B$。

联结词 $\lif$ 是\emph{真值函数性的}，即 $!A \lif !B$ 的真值——$\True$ 或 $\False$——由 $!A$ 和 $!B$ 的真值决定：$!A \lif !B$ 为真，当且仅当 $!A$ 为假或 $!B$ 为真，否则为假。相对于一个真值赋值 $\pAssign v$，我们定义 $\pSat{v}{!A \lif !B}$ 当且仅当 $\pSat/{v}{!A}$ 或 $\pSat{v}{!B}$。具有这种语义的联结词 $\lif$ 被称为\emph{实质条件句}。

根据这一定义，可以得出许多基本的逻辑事实。首先，对于实质条件句，演绎定理成立：
\begin{align}
  \text{若 } \Gamma, !A \Entails !B \text{ 则 } \Gamma \Entails !A \lif !B
\end{align}
它是真值函数性的：$!A \lif !B$ 与 $\lnot !A \lor !B$ 等值：
\begin{align}
  !A \lif !B & \Entails \lnot !A \lor !B\\
  \lnot !A \lor !B & \Entails !A \lif !B
  \intertext{一个实质条件句是其后件及其前件否定的语义后承：}
  !B & \Entails !A \lif !B\\
  \lnot !A & \Entails !A \lif !B
  \intertext{一个假的实质条件句等值于其前件与其后件否定的合取：如果 $!A \lif !B$ 为假，则 $!A \land \lnot !B$ 为真，反之亦然：}
  \lnot(!A \lif !B) & \Entails !A \land \lnot !B\\
  !A \land \lnot !B & \Entails \lnot(!A \lif !B)
  \intertext{实质条件句支持肯定前件式：}
  !A, !A \lif !B & \Entails !B
  \intertext{实质条件句可聚合：}
  !A \lif !B,  !A \lif !C & \Entails !A \lif (!B \land !C)
  \intertext{我们总是可以强化前件，即条件句是\emph{单调的}：}
  !A \lif !B & \Entails (!A \land !C) \lif !B
  \intertext{实质条件句是可传递的，即连锁推理有效：}
  !A \lif !B, !B \lif !C & \Entails !A \lif !C
  \intertext{实质条件句与其逆否命题等值：}
  !A \lif !B & \Entails \lnot !B \lif \lnot !A\\
  \lnot !B \lif \lnot !A & \Entails !A \lif !B
\end{align}

这些在数学推理中都是有用且无疑问的推理。然而，哲学和语言学文献中充斥着针对非数学语境中等价推理的所谓反例。这些反例表明，在符号化英语“如果……那么……”陈述时，实质条件句 $\lif$ 并不是——或至少不总是——合适的联结词。

\end{document}